% §16.6 What's Next
%
% Cross-reference resolution status:
%
% DEFINED here:
%   \label{sec:ch16-whats-next}                 -- this section
%
% RESOLVED (within chapter 16):
%   \Cref{sec:varying-rates}                    -- §16.2
%   \Cref{sec:correlated-noise}                 -- §16.3
%   \Cref{sec:adversarial-inputs}               -- §16.4
%   \Cref{sec:where-framework-helps}            -- §16.5
%   \Cref{sec:bib-notes-ch16}                   -- §16.7 (Task 8)
%
% RESOLVED (chapter labels for forward references):
%   \Cref{ch:not-bernoulli}                     -- chapter 16 (this chapter)
%   Chapter 17 (cipher maps and trapdoor computing) -- Plan 0 stub
%   Chapter 18 (encrypted search)               -- Plan 0 stub
%
% §16.6 closes the negative-space catalogue, signposts the
% cryptographic applications of chapters 17-18, and forward-references
% the bibliographic notes for the analytical alternatives the chapter
% has named.

\section{What's Next}
\label{sec:ch16-whats-next}

The chapter has completed the boundary cataloguing the framework
called for at the outset of Part~V. The negative space around the
two axioms is now mapped in three sections, each named by the
mechanism that pushes a system off the iid regime: varying rates
over time (\Cref{sec:varying-rates}), correlated noise across
queries (\Cref{sec:correlated-noise}), and adversarial inputs
(\Cref{sec:adversarial-inputs}). \Cref{sec:where-framework-helps}
recovered the per-query content of the framework that survives
each of these failures, leaving the practitioner with a working
rule rather than a binary verdict: the geometric and channel-matrix
machinery applies per-query; the multi-query bounds require
replacement by their non-iid counterparts. With the boundary
mapped, the book's remaining chapters can take up two threads that
the boundary work has only signposted.

\subsection*{Closing the negative space}

\Cref{ch:not-bernoulli} has been a chapter of qualifications: a
catalogue of what the framework does not cover and what to do
instead. Its placement between the constructive culmination of
Part~IV (the BHF and BPHF) and the cryptographic applications of
Part~VI is structural. Part~IV sat inside the framework's domain
of validity; Part~VI extends the framework into settings where
clean iid analysis can no longer be relied on. The boundary
chapter is the explicit bridge between the two. Without it, the
applications would inherit the iid assumptions implicitly, and
the guarantees they advertise would lean on hidden conditions.

The chapter has also closed five deferrals seeded across the
preceding parts: the time-varying-rate question from chapter~1
on sampling, the correlated-witness question from chapter~3 and
reopened in chapter~1's exercises, the adaptive-input question
from chapter~5 on one-sided sets, the adversarial-stream question
from chapter~10 on sketch composition, and the algebra-types
question from the chapter~13 exercises. Each deferral now has a
section assigned to it, with the violated axiom named and the
analytical replacement identified. No further deferrals to this
chapter remain open.

\subsection*{Forward to chapters 17 and 18}

The two chapters after this one apply the Bernoulli framework
under cryptographic constraints, and both are direct heirs to the
boundary work the present chapter has done. Chapter~17 develops
cipher maps and trapdoor computing as the algebraic framework
that recasts the BHF and BPHF architecture for settings where the
underlying perfect hash is replaced by a trapdoor function: the
forward direction is public and computable by anyone, the reverse
direction is private and computable only with a secret key.
Chapter~18 specializes the cipher-map machinery to encrypted
search, showing how the FPR and FNR guarantees of Part~IV carry
over into the cryptographic setting and what additional structure
(oblivious access patterns, side-channel considerations) the
cryptographic context demands.

Both chapters live in territory that the present chapter has
flagged as non-trivial. Encrypted-search queries are adversarial
by construction: the server is curious and the queries are chosen
by an attacker who has access to the encrypted index and may
adapt to responses. Chapter~17 names the worst-case analysis
tools as the replacement for iid expectations, and chapter~18
shows how cryptographic hash families close the side channels
\Cref{sec:adversarial-inputs} identified.

\subsection*{The analytical alternatives}

The chapter has named, but not bibliographically anchored, the
analytical alternatives the non-iid regime requires. The Azuma-
Hoeffding inequality of \Cref{sec:varying-rates}, McDiarmid's
inequality of \Cref{sec:correlated-noise}, and the worst-case and
cryptographic-hash-family tools of \Cref{sec:adversarial-inputs}
all sit in literatures whose entry points the bibliographic notes
of \Cref{sec:bib-notes-ch16} attach to the chapter. The notes are
brief by design; the chapter's purpose has been to name the right
tool for each failure mode and to give the reader a working
intuition for when each applies. The notes point the reader to
the foundational papers and to a standard textbook reference so
that the practitioner can pick up the trail without having to
triangulate from the chapter's prose alone.
